\chapter{1988年全国卷（理）}

\section{单选题}

\begin{problem}

\(\left(\frac{1 - \i}{1 + \i}\right)^2\) 的值等于（\(\quad\)）

\choices
    {\(1\)}
    {\(-1\)}
    {\(\i\)}
    {\(-\i\)}
\end{problem}

\begin{answer}
B
\end{answer}

\begin{solution}
\[
\frac{1-\i}{1+\i}=\frac{(1-\i)^2}{(1+\i)(1-\i)}=\frac{-2\i}{2}=-\i
\]
因此
\[
\left(\frac{1-\i}{1+\i}\right)^2=(-\i)^2=-1
\]
故选 B.
\end{solution}


\begin{problem}

设圆 \(M\) 的方程为 \((x - 3)^2 + (y - 2)^2 = 2\)，直线 \(L\) 的方程为 \(x + y - 3 = 0\)，点 \(P\) 的坐标为 \((2,1)\). 那么（\(\quad\)）

\choices
    {点 \(P\) 在直线 \(L\) 上，但不在圆 \(M\) 上}
    {点 \(P\) 在圆 \(M\) 上，但不在直线 \(L\) 上}
    {点 \(P\) 既在圆 \(M\) 上，又在直线 \(L\) 上}
    {点 \(P\) 既不在直线 \(L\) 上，也不在圆 \(M\) 上}
\end{problem}

\begin{answer}
C
\end{answer}

\begin{solution}
将点 \(P(2,1)\) 代入直线 \(L\) 的方程，得 \(2+1-3=0\)，所以 \(P\in L\). 又
\[
(2-3)^2+(1-2)^2=1+1=2
\]
满足圆 \(M\) 的方程，所以 \(P\in M\). 故选 C.
\end{solution}

\begin{problem}

集合 \(\{1,2,3\}\) 的子集共有（\(\quad\)）

\choices
    {\(7\) 个}
    {\(8\) 个}
    {\(6\) 个}
    {\(5\) 个}
\end{problem}

\begin{answer}
B
\end{answer}

\begin{solution}
集合中有 \(3\) 个元素，每个元素对于一个子集都有“选入”或“不选入”两种可能，且这些选择彼此独立.因此子集总数为
\[
2^3=8
\]
其中包括空集和集合本身，故选 B.
\end{solution}

\begin{problem}

已知双曲线方程 \(\frac{x^2}{20} - \frac{y^2}{5} = 1\)，那么它的焦距是（\(\quad\)）

\choices
    {\(10\)}
    {\(5\)}
    {\(\sqrt{15}\)}
    {\(2\sqrt{15}\)}
\end{problem}

\begin{answer}
A
\end{answer}

\begin{solution}
双曲线标准方程中的 \(a^2=20\)，\(b^2=5\)，所以焦半距 \(c\) 满足
\(c^2=a^2+b^2=20+5=25\)，\(c=5\).
焦距是两焦点之间的距离 \(2c=10\)，故选 A.
\end{solution}

\begin{problem}

在 \((x - \sqrt{3})^10\) 的展开式中，\(x^6\) 的系数是（\(\quad\)）

\choices
    {\(-27\mathrm C_{10}^6\)}
    {\(27\mathrm C_{10}^4\)}
    {\(-9\mathrm C_{10}^6\)}
    {\(9\mathrm C_{10}^4\)}
\end{problem}

\begin{answer}
D
\end{answer}

\begin{solution}
二项式展开的一般项中，若取 \(x^6\)，就必须从另外的 \(4\) 项中取 \(-\sqrt{3}\). 因而 \(x^6\) 的系数为
\[
\mathrm C_{10}^{4}(-\sqrt{3})^4=\mathrm C_{10}^{4}\cdot9=9\mathrm C_{10}^{4}
\]
故选 D.
\end{solution}

\begin{problem}

函数 \(y=\cos^4 x-\sin^4 x\) 的最小正周期是（\(\quad\)）

\choices
    {\(\pi\)}
    {\(2\pi\)}
    {\(\frac{\pi}{2}\)}
    {\(4\pi\)}
\end{problem}

\begin{answer}
A
\end{answer}

\begin{solution}
\[
y=\cos^4x-\sin^4x
=(\cos^2x-\sin^2x)(\cos^2x+\sin^2x)
=\cos 2x
\]
函数\(\cos 2x\)的周期为\(\pi\)，且\(\pi\)是最小的正周期，因此最小正周期为\(\pi\)，故选 A.
\end{solution}

\begin{problem}

方程 \(4\cos^2 x - 4\sqrt{3}\cos x + 3 = 0\) 的解集是（\(\quad\)）

\choices
    {\(\left\{x\mid x = k\pi +(-1)^k\cdot \frac{\pi}{6},k\in \Z\right\}\)}
    {\(\left\{x\mid x = k\pi +(-1)^k\cdot \frac{\pi}{3},k\in \Z\right\}\)}
    {\(\left\{x\mid x = 2k\pi \pm \frac{\pi}{6},k\in \Z\right\}\)}
    {\(\left\{x\mid x = 2k\pi \pm \frac{\pi}{3},k\in \Z\right\}\)}
\end{problem}

\begin{answer}
C
\end{answer}

\begin{solution}
\[
4\cos^2x-4\sqrt{3}\cos x+3
=(2\cos x-\sqrt{3})^2=0
\]
所以\(\cos x=\frac{\sqrt{3}}{2}\)，从而
\(x=2k\pi\pm\frac{\pi}{6}\)，\(k\in\Z\).
这正是选项 C，故选 C.
\end{solution}

\begin{problem}

极坐标方程 \(\rho = \frac{4}{3 - 2\cos\theta}\) 所表示的曲线是（\(\quad\)）

\choices
    {圆}
    {双曲线右支}
    {抛物线}
    {椭圆}
\end{problem}

\begin{answer}
D
\end{answer}

\begin{solution}
由极坐标与直角坐标的关系\(x=\rho\cos\theta\)，原方程化为
\[
3\rho-2\rho\cos\theta=4
\quad\Longleftrightarrow\quad
\rho=\frac{2}{3}(x+2)
\]
由于\(3-2\cos\theta\ge1\)，有\(\rho>0\)，从而由上式知\(x+2>0\). 因此\(\rho\)是点到原点的距离，\(x+2=\lvert x+2\rvert\)是点到直线\(x=-2\)的距离，故该曲线的离心率为\(\frac{2}{3}<1\)，表示椭圆，故选 D.
\end{solution}

\begin{problem}

如图，正四棱台中，\(A'D'\) 所在的直线与 \(BB'\) 所在的直线是（\(\quad\)）

\bitmapfigure[width=0.46\linewidth]{img/1988/national_science/q9_fig1.png}

\choices
    {相交直线}
    {平行直线}
    {不互相垂直的异面直线}
    {互相垂直的异面直线}
\end{problem}

\begin{answer}
C
\end{answer}

\begin{solution}
设正四棱台的下底面为\(z=0\)，上底面为\(z=h\).取下底面中心为原点，并令下底面四个顶点中
\(A=(-1,-1,0)\)，\(B=(1,-1,0)\)，\(D=(-1,1,0)\).
设上底面边长与下底面边长之比为\(\lambda\)，则\(0<\lambda<1\)，可取
\(A'=(-\lambda,-\lambda,h)\)，\(B'=(\lambda,-\lambda,h)\)，\(D'=(-\lambda,\lambda,h)\).
于是两条直线的方向向量可分别取为
\(\overrightarrow{A'D'}=(0,2\lambda,0)\)，\(\overrightarrow{BB'}=(\lambda-1,1-\lambda,h)\).
它们的内积为\(2\lambda(1-\lambda)\ne0\)，所以两直线不垂直.直线\(A'D'\)完全在上底面内，而直线\(BB'\)与上底面只有交点\(B'\)，且\(B'\notin A'D'\)，所以两直线不相交；方向向量也不成比例，故不平行.因此它们是不互相垂直的异面直线，故选 C.
\end{solution}

\begin{problem}

\(\tan \left(\arctan \frac{1}{5}+\arctan 3\right)\) 的值等于（\(\quad\)）

\choices
    {\(4\)}
    {\(\frac{1}{2}\)}
    {\(\frac{1}{8}\)}
    {\(8\)}
\end{problem}

\begin{answer}
D
\end{answer}

\begin{solution}
利用正切的和角公式，得
\[
\tan\left(\arctan\frac{1}{5}+\arctan 3\right)
=\frac{\frac{1}{5}+3}{1-\frac{1}{5}\cdot3}
=\frac{\frac{16}{5}}{\frac{2}{5}}=8
\]
分母\(1-\frac{3}{5}\ne0\)，故运算合法，答案为\(8\)，故选 D.
\end{solution}

\begin{problem}

设命题甲：\(\triangle ABC\) 的一个内角为 \(60^{\circ}\). 命题乙：\(\triangle ABC\) 的三内角的度数成等差数列. 那么（\(\quad\)）

\choices
    {甲是乙的充分条件，但不是必要条件}
    {甲是乙的必要条件，但不是充分条件}
    {甲是乙的充要条件}
    {甲不是乙的充分条件，也不是乙的必要条件}
\end{problem}

\begin{answer}
C
\end{answer}

\begin{solution}
设三角形\(ABC\)的三个内角按从小到大排列为\(\alpha\le\beta\le\gamma\).若命题甲成立，某一个内角为\(60^{\circ}\).若这个角是中间角，则\(\beta=60^{\circ}\)，而\(\alpha+\gamma=120^{\circ}=2\beta\)，所以三个角成等差数列.

若\(60^{\circ}\)是最小角，则\(\beta\)，\(\gamma\ge60^{\circ}\)且\(\beta+\gamma=120^{\circ}\)，从而\(\beta=\gamma=60^{\circ}\)；若\(60^{\circ}\)是最大角，同理也得到三个角均为\(60^{\circ}\).因此甲是乙的充分条件.

反过来，若三个内角成等差数列，则中间角为三角形内角和的三分之一，即\(\beta=\frac{180^{\circ}}{3}=60^{\circ}\)，所以乙是甲的充分条件.故甲是乙的充要条件，选 C.
\end{solution}

\begin{problem}

在复平面内，若复数 \(z\) 满足 \(|z + 1| = |z - \i|\)，则 \(z\) 所对应的点 \(Z\) 的集合构成的图形是（\(\quad\)）

\choices
    {圆}
    {直线}
    {椭圆}
    {双曲线}
\end{problem}

\begin{answer}
B
\end{answer}

\begin{solution}
设\(z=x+\i y\)，其中\(x\)，\(y\in\R\).则
\(|z+1|^2=(x+1)^2+y^2\)，\(|z-\i|^2=x^2+(y-1)^2\).
由题设两式相等，得
\[
(x+1)^2+y^2=x^2+(y-1)^2
\quad\Longleftrightarrow\quad
x+y=0
\]
因此点\(Z\)的轨迹是直线\(x+y=0\)，故选 B.
\end{solution}

\begin{problem}

如果曲线 \(x^{2} - y^{2} - 2x - 2y - 1 = 0\) 经过平移坐标轴后的新方程为 \(x^{\prime 2} - y^{\prime 2} = 1\)，那么新坐标系的原点在原坐标系中的坐标为（\(\quad\)）

\choices
    {\((1,1)\)}
    {\((-1, - 1)\)}
    {\((-1, 1)\)}
    {\((1, - 1)\)}
\end{problem}

\begin{answer}
D
\end{answer}

\begin{solution}
将原方程配方，得
\[
x^2-y^2-2x-2y-1=0
\quad\Longleftrightarrow\quad
(x-1)^2-(y+1)^2=1
\]
令新坐标为\(x'=x-1\)，\(y'=y+1\)，便得到题设的新方程\(x'^2-y'^2=1\).新坐标系原点满足\(x'=0\)、\(y'=0\)，故其在原坐标系中的坐标为\((1,-1)\)，选 D.
\end{solution}

\begin{problem}

假设在 \(200\) 件产品中有 \(3\) 件次品，现在从中任意抽取 \(5\) 件，其中至少有 \(2\) 件次品的抽法有（\(\quad\)）

\choices
    {\(\mathrm C_3^2\mathrm C_{197}^3\) 种}
    {\(\mathrm C_3^2\mathrm C_{197}^3 + \mathrm C_3^3\mathrm C_{197}^2\) 种}
    {\(\mathrm C_{200}^5 - \mathrm C_{197}^5\) 种}
    {\(\mathrm C_{200}^5 - \mathrm C_3^1\mathrm C_{197}^4\) 种}
\end{problem}

\begin{answer}
B
\end{answer}

\begin{solution}
至少有\(2\)件次品的情况只有两类：恰有\(2\)件次品，或恰有\(3\)件次品.分别计数得
\[
\mathrm C_3^2\mathrm C_{197}^3
\quad\text{和}\quad
\mathrm C_3^3\mathrm C_{197}^2
\]
两类情况互不相交，所以抽法总数为
\[
\mathrm C_3^2\mathrm C_{197}^3+\mathrm C_3^3\mathrm C_{197}^2
\]
故选 B.
\end{solution}

\begin{problem}

已知二面角 \(\alpha-AB-\beta\) 的平面角是锐角，\(C\) 是平面 \(\alpha\) 内一点（它不在棱 \(AB\) 上），点 \(D\) 是点 \(C\) 在面 \(\beta\) 上的射影，点 \(E\) 是棱 \(AB\) 上满足 \(\angle CEB\) 为锐角的任一点，那么（\(\quad\)）

\choices
    {\(\angle CEB > \angle DEB\)}
    {\(\angle CEB = \angle DEB\)}
    {\(\angle CEB < \angle DEB\)}
    {\(\angle CEB\) 与 \(\angle DEB\) 的大小关系不能确定}
\end{problem}

\begin{answer}
A
\end{answer}

\begin{solution}
设二面角的平面角为\(\theta\)，则\(0<\theta<\frac{\pi}{2}\).在过\(E\)且垂直于\(AB\)的截面内，取\(EB\)方向为单位向量\(\bs e\)，取平面\(\alpha\)、\(\beta\)在该截面内的方向单位向量分别为\(\bs u\)、\(\bs v\)，于是\(\bs u\cdot\bs v=\cos\theta\).

由于\(C\)不在\(AB\)上，可写成
\(\overrightarrow{EC}=x\bs e+r\bs u\)，\(r>0\).
条件\(\angle CEB\)为锐角说明\(x>0\).点\(D\)是\(C\)在平面\(\beta\)上的射影，正交投影保留沿\(AB\)的分量并把\(\bs u\)投影为\(\cos\theta\bs v\)，所以
\[
\overrightarrow{ED}=x\bs e+r\cos\theta\bs v
\]
因此\(\angle DEB\)也为锐角，并且
\(\tan\angle CEB=\frac{r}{x}\)，\(\tan\angle DEB=\frac{r\cos\theta}{x}\).
因为\(0<\cos\theta<1\)，所以\(\tan\angle DEB<\tan\angle CEB\)，从而\(\angle CEB>\angle DEB\)，故选 A.
\end{solution}

\section{填空题}

\begin{problem}

求复数 \(\sqrt{3} - \i\) 的模和辐角的主值，分别为\fillinblank{}，\fillinblank{}.
\end{problem}

\begin{answer}
\(2\)，\(\frac{11\pi}{6}\)
\end{answer}

\begin{solution}
复数\(z=\sqrt{3}-\i\)对应的点为\((\sqrt{3},-1)\)，所以
\[
|z|=\sqrt{(\sqrt{3})^2+(-1)^2}=\sqrt4=2
\]
该点在第四象限，且参考角为\(\frac{\pi}{6}\).按本题所用的辐角主值范围\([0,2\pi)\)，主值为
\[
\arg z=2\pi-\frac{\pi}{6}=\frac{11\pi}{6}
\]
\end{solution}

\begin{problem}

解方程：\(9^{-x} - 2 \cdot 3^{1 - x} = 27\)，解为\fillinblank{}.
\end{problem}

\begin{answer}
\(-2\)
\end{answer}

\begin{solution}
令\(t=3^{-x}>0\)，则\(9^{-x}=3^{-2x}=t^2\)，且\(3^{1-x}=3t\).原方程化为
\[
t^2-6t=27
\quad\Longleftrightarrow\quad
t^2-6t-27=0
\quad\Longleftrightarrow\quad
(t-9)(t+3)=0
\]
因为\(t>0\)，只能取\(t=9=3^2\)，所以\(3^{-x}=3^2\)，解得\(x=-2\).
\end{solution}

\begin{problem}

已知 \(\sin \theta = -\frac{3}{5}\)，\(3\pi < \theta < \frac{7\pi}{2}\)，求 \(\tan \frac{\theta}{2}\) 的值为\fillinblank{}.
\end{problem}

\begin{answer}
\(-3\)
\end{answer}

\begin{solution}
由\(3\pi<\theta<\frac{7\pi}{2}\)知\(\theta\)在第三象限，因此\(\cos\theta<0\).由
\(\sin^2\theta+\cos^2\theta=1\)，\(\sin\theta=-\frac35\)
得\(\cos\theta=-\frac45\).又\(\frac{\theta}{2}\in(\frac{3\pi}{2},\frac{7\pi}{4})\)，半角正切公式给出
\[
\tan\frac{\theta}{2}=\frac{\sin\theta}{1+\cos\theta}
=\frac{-\frac35}{1-\frac45}=-3
\]
\end{solution}

\begin{problem}

如图，四棱锥 \(S-ABCD\) 的底面是边长为 \(1\) 的正方形，侧棱 \(SB\) 垂直于底面，并且 \(SB=\sqrt{3}\)，用 \(\alpha\) 表示 \(\angle ASD\)，求 \(\sin \alpha\) 的值为\fillinblank{}.

\bitmapfigure[width=0.32\linewidth]{img/1988/national_science/q19_fig1.png}
\end{problem}

\begin{answer}
\(\frac{\sqrt{5}}{5}\)
\end{answer}

\begin{solution}
取底面所在平面为\(z=0\)，令\(B=(0,0,0)\)，\(A=(1,0,0)\)，\(C=(0,1,0)\)，\(D=(1,1,0)\).因为\(SB\perp\)底面且\(SB=\sqrt3\)，可取\(S=(0,0,\sqrt3)\).于是
\(\overrightarrow{SA}=(1,0,-\sqrt3)\)，\(\overrightarrow{SD}=(1,1,-\sqrt3)\).
所以
\(|\overrightarrow{SA}|=2\)，\(|\overrightarrow{SD}|=\sqrt5\)，\(\overrightarrow{SA}\cdot\overrightarrow{SD}=1+3=4\).
故
\(\cos\alpha=\frac{4}{2\sqrt5}=\frac{2}{\sqrt5}\)，\(\sin\alpha=\sqrt{1-\frac45}=\frac{\sqrt5}{5}\).
\end{solution}

\begin{problem}

已知等比数列 \(\{a_{n}\}\) 的公比 \(q > 1\)，并且 \(a_{1} = b\)（\(b \neq 0\)）. 求 \(\displaystyle\lim_{n \to \infty} \frac{a_{1} + a_{2} + a_{3} + \cdots + a_{n}}{a_{6} + a_{7} + a_{8} + \cdots + a_{n}}=\)\fillinblank{}.
\end{problem}

\begin{answer}
\(1\)
\end{answer}

\begin{solution}
等比数列的通项为\(a_k=bq^{k-1}\).当\(n\ge6\)时，
\(\sum_{k=1}^{n}a_k=b\frac{q^n-1}{q-1}\)，\(\sum_{k=6}^{n}a_k=bq^5\frac{q^{n-5}-1}{q-1}=b\frac{q^n-q^5}{q-1}\).
因此
\[
\frac{a_1+a_2+\cdots+a_n}{a_6+a_7+\cdots+a_n}
=\frac{q^n-1}{q^n-q^5}
\]
由于\(q>1\)，有\(q^n\to\infty\)，故
\[
\lim_{n\to\infty}\frac{q^n-1}{q^n-q^5}=1
\]
\end{solution}

\section{解答题}

\begin{problem}

已知 \(\tan x = a\)，求 \(\frac{3\sin x + \sin 3x}{3\cos x + \cos 3x}\) 的值.
\end{problem}

\begin{solution}
由\(\tan x=a\)知\(\cos x\ne0\)，并且\(\sin x=a\cos x\).利用三倍角公式，写成\(a\)的形式为
\[
\sin3x=\sin x(3-4\sin^2x)
=\frac{a\cos x(3-a^2)}{1+a^2}
\]
\[
\cos3x=\cos x(1-4\sin^2x)
=\frac{\cos x(1-3a^2)}{1+a^2}
\]
于是
\[
3\sin x+\sin3x
=\frac{2a\cos x(a^2+3)}{1+a^2}
\]
\[
3\cos x+\cos3x
=\frac{4\cos x}{1+a^2}\ne0
\]
所以
\[
\frac{3\sin x+\sin3x}{3\cos x+\cos3x}
=\frac{a(a^2+3)}{2}
\]
\end{solution}

\begin{problem}

如图，正三棱锥 \(S-ABC\) 的侧面是边长为 \(a\) 的正三角形，\(D\) 是 \(SA\) 的中点，\(E\) 是 \(BC\) 的中点，求 \(\triangle SDE\) 绕直线 \(SE\) 旋转一周所得的旋转体的体积.
\end{problem}

\begin{solution}
连接\(AE\).由于\(\triangle SBC\)为边长为\(a\)的正三角形，\(E\)是\(BC\)的中点，所以
\[
SE=\frac{\sqrt3}{2}a
\]
又\(D\)是\(SA\)的中点，故\(SD=\frac a2\).利用\(S\)点处三条棱两两成\(60^{\circ}\)，有
\[
\overrightarrow{DE}=\overrightarrow{SE}-\overrightarrow{SD}
=\frac12\left(\overrightarrow{SB}+\overrightarrow{SC}-\overrightarrow{SA}\right)
\]
又\(\lvert\overrightarrow{SA}\rvert=\lvert\overrightarrow{SB}\rvert=\lvert\overrightarrow{SC}\rvert=a\)，且任意两条向量的夹角为\(60^{\circ}\)，所以
\[
DE^2=\frac14\left\lvert\overrightarrow{SB}+\overrightarrow{SC}-\overrightarrow{SA}\right\rvert^2
=\frac14\left(3a^2+2\cdot\frac{a^2}{2}-2\cdot\frac{a^2}{2}-2\cdot\frac{a^2}{2}\right)
=\frac{a^2}{2}
\]
在\(\triangle SDE\)中，
\[
\cos\angle DSE
=\frac{SD^2+SE^2-DE^2}{2\cdot SD\cdot SE}
=\frac{1}{\sqrt3}
\]
设\(F\)为\(D\)到\(SE\)的垂足，则
\(SF=SD\cos\angle DSE=\frac{\sqrt3}{6}a\)，\(EF=SE-SF=\frac{\sqrt3}{3}a\).
此外，
\[
DF=SD\sin\angle DSE
=\frac a2\sqrt{1-\frac13}=\frac{\sqrt6}{6}a
\]
旋转后所得几何体是以\(DF\)为底面半径、分别以\(SF\)和\(EF\)为高的两个圆锥的组合，因此
\[
V=\frac13\pi DF^2\cdot SF+\frac13\pi DF^2\cdot EF
=\frac13\pi\left(\frac{\sqrt6}{6}a\right)^2
\left(\frac{\sqrt3}{6}a+\frac{\sqrt3}{3}a\right)
=\frac{\sqrt3}{36}\pi a^3
\]
\end{solution}

\begin{problem}

设 \(a > 0\)，\(a \neq 1\)，\(t > 0\)，比较 \(\frac{1}{2} \log_a t\) 与 \(\log_a \frac{t + 1}{2}\) 的大小，并证明你的结论.
\end{problem}

\begin{solution}
考虑两者之差：
\[
\frac12\log_a t-\log_a\frac{t+1}{2}
=\frac{1}{\ln a}\ln\frac{2\sqrt t}{t+1}
\]
由算术几何平均不等式，\(t+1\ge2\sqrt t\)，且当且仅当\(t=1\)取等号.因此当\(t=1\)时两者相等；当\(t\ne1\)时，
\[
0<\frac{2\sqrt t}{t+1}<1
\]
从而\(\ln\frac{2\sqrt t}{t+1}<0\).若\(a>1\)，则\(\ln a>0\)，所以
\[
\frac12\log_a t<\log_a\frac{t+1}{2}
\]
若\(0<a<1\)，则\(\ln a<0\)，不等号反向，所以
\[
\frac12\log_a t>\log_a\frac{t+1}{2}
\]
\end{solution}

\begin{problem}

给定实数 \(a\)，且 \(a \neq 0\)，\(a \neq 1\)，设函数 \(y = \frac{x - 1}{ax - 1}\)（\(x \in \R\)，且 \(x \neq \frac{1}{a}\)）. 证明：
\begin{enumerate}
\item 经过这个函数图象上任意两个不同的点的直线不平行于 \(x\) 轴；
\item 这个函数的图象关于直线 \(y = x\) 成轴对称图形.
\end{enumerate}
\end{problem}

\begin{solution}
\begin{enumerate}
\item 设图象上两个不同的点为\(M_1(x_1,y_1)\)、\(M_2(x_2,y_2)\)，则\(x_1\ne x_2\)，且\(ax_1-1\ne0\)、\(ax_2-1\ne0\).由\(y_i=\frac{x_i-1}{ax_i-1}\)，有
\[
y_1-y_2
=\frac{(a-1)(x_1-x_2)}{(ax_1-1)(ax_2-1)}\ne0
\]
因此\(\frac{y_1-y_2}{x_1-x_2}\ne0\)，经过这两个点的直线不平行于\(x\)轴.

\item 设\(P(x,y)\)在函数图象上，则
\[
y=\frac{x-1}{ax-1}
\quad\Longleftrightarrow\quad
axy-y-x+1=0
\]
这个方程关于\(x\)、\(y\)对称.还需验证交换后分母不为零：若\(ay-1=0\)，代入上式可得\(1-\frac1a=0\)，即\(a=1\)，与题设矛盾，所以\(ay-1\ne0\).于是
\[
x=\frac{y-1}{ay-1}
\]
说明点\((y,x)\)也在函数图象上.因此图象关于直线\(y=x\)成轴对称图形.
\end{enumerate}
\end{solution}

\begin{problem}

直线 \(L\) 的方程为 \(x = -\frac{p}{2}\)，其中 \(p > 0\). 椭圆的中心为 \(D\left(2+\frac{p}{2},0\right)\)，焦点在 \(x\) 轴上，长半轴长为 \(2\)，短半轴长为 \(1\)，它的一个顶点为 \(A\left(\frac{p}{2},0\right)\). 问 \(p\) 在哪个范围内取值时，椭圆上有四个不同的点，它们中每一个点到点 \(A\) 的距离等于该点到直线 \(L\) 的距离.
\end{problem}

\begin{solution}
椭圆的方程为
\[
\frac{\left(x-\left(2+\frac p2\right)\right)^2}{4}+y^2=1
\]
点\((x,y)\)到\(A\left(\frac p2,0\right)\)的距离等于它到直线\(L:x=-\frac p2\)的距离，当且仅当
\[
\left(x-\frac p2\right)^2+y^2=\left(x+\frac p2\right)^2
\quad\Longleftrightarrow\quad
y^2=2px
\]
与椭圆联立，得到
\(x^2+(7p-4)x+\frac{p^2}{4}+2p=0\)，\(y^2=2px\).
要得到四个不同的点，关于\(x\)的方程必须有两个不相等的正根：每个正根都对应两个不同的\(y=\pm\sqrt{2px}\).其判别式为
\[
\Delta=(7p-4)^2-4\left(\frac{p^2}{4}+2p\right)
=16(3p-1)(p-1)>0
\]
又因为两根之积\(\frac{p^2}{4}+2p>0\)，两根之和为\(4-7p\)，所以两根均为正还需\(4-7p>0\).结合\(p>0\)可得
\[
0<p<\frac13
\]
在此范围内\(\Delta>0\)，且两根均为正，每个根确实给出\(y=\pm\sqrt{2px}\)两个不同的点，故恰有四个不同的点.所求范围为\(0<p<\frac13\).
\end{solution}
